\usepackage{graphicx} % Для \includegraphics
\usepackage{rotating}
\usepackage{adjustbox}
\RequirePackage{float}
\RequirePackage{wrapfig}
\RequirePackage{tikzscale}
\RequirePackage{rotating} % For \rotatebox command
\RequirePackage{calc}  

% Создание шапки титульной страницы
\renewcommand{\titlepageheader}[2]
{
	\begin{wrapfigure}[7]{l}{0.14\linewidth}
		\vspace{3.4mm}
		\hspace{-8mm}
		\includegraphics[width=0.89\linewidth]{bmstu-logo}
	\end{wrapfigure}
	
	{
		\singlespacing \small
		Министерство науки и высшего образования Российской Федерации \\
		Федеральное государственное автономное образовательное учреждение \\
		высшего образования \\
		<<Московский государственный технический университет \\
		имени Н.~Э.~Баумана \\
		(национальный исследовательский университет)>> \\
		(МГТУ им. Н.~Э.~Баумана) \\
	}
	
	\vspace{-4.2mm}
	\vhrulefill{0.9mm} \\
	\vspace{-7mm}
	\vhrulefill{0.2mm} \\
	\vspace{2.8mm}
	
	{
		\small
		ФАКУЛЬТЕТ \longunderline{<<#1>>} \\
		\vspace{3.3mm}
		КАФЕДРА \longunderline{<<#2>>} \\
	}
}

% Установка заголовков РПЗ
\renewcommand{\titlepagenotetitle}[2]
{
	{
		\LARGE \bfseries
		РАСЧЕТНО-ПОЯСНИТЕЛЬНАЯ ЗАПИСКА \\
	}
	\vspace{5mm}
	{
		\Large \itshape
		#1 \\
		\vspace{5mm}
		НА ТЕМУ: \\
		\vspace{5mm}
		<<#2>> \\
	}
}

% Создание страницы определений
\renewenvironment{definitions}
{
	\chapter*{ОПРЕДЕЛЕНИЯ}
	\addcontentsline{toc}{chapter}{ОПРЕДЕЛЕНИЯ}
	
	В настоящей расчетно-пояснительной записке применяют следующие термины с соответствующими определениями.
	
	\begin{description}[leftmargin=0pt]
	}
	{
	\end{description}
}

% Создание страницы обозначений и сокращений
\renewenvironment{abbreviations}
{
	\chapter*{ОБОЗНАЧЕНИЯ И СОКРАЩЕНИЯ}
	\addcontentsline{toc}{chapter}{ОБОЗНАЧЕНИЯ И СОКРАЩЕНИЯ}
	
	В настоящей расчетно-пояснительной записке применяют следующие сокращения и обозначения.
	
	\begin{description}[leftmargin=0pt]
	}
	{
	\end{description}
}

\RequirePackage{enumitem}
\newcommand{\abbreviation}[2]
{
	\item \noindent #1 --- #2
}

\newcommand{\ssr}[1]
{
	\chapter{#1}
	\addcontentsline{toc}{chapter}{#1}  
}

\newcommand{\ssrno}[1]
{
	\chapter*{#1}
	\addcontentsline{toc}{chapter}{#1}  
}


% Переопределение формата даты обращения
\DeclareFieldFormat{urldate}
{%
	(дата~обращения:%
	\addspace
	\iffieldundef{urlday}{}{\ifnumless{\thefield{urlday}}{10}{0\thefield{urlday}}{\thefield{urlday}}%
		\adddot
		\iffieldundef{urlmonth}{}{\ifnumless{\thefield{urlmonth}}{10}{0\thefield{urlmonth}}{\thefield{urlmonth}}%
			\adddot
			\thefield{urlyear}}})
}

% Improved command for creating a figure
\renewcommand{\includeimage}[6]% Now has 6 mandatory arguments
{
	\ifthenelse{\equal{#2}{f}}
	{
		\ifthenelse{\equal{#6}{y}}% Check the rotation flag
		{
			% Use a standard figure environment and rotate contents with adjustbox
			\begin{figure}[#3] % The placement specifier #3 is now used
				\centering
				\begin{adjustbox}{addcode={\begin{minipage}{\width}}{\caption{#5}\label{img:#1}\end{minipage}},rotate=90,center}
					\includegraphics[width=#4]{inc/img/#1}%
				\end{adjustbox}
			\end{figure}
		}% Rotate if 'y'
		{
			% Standard figure if no rotation is needed
			\begin{figure}[#3]
				\centering
				\includegraphics[width=#4]{inc/img/#1}
				\caption{#5}
				\label{img:#1}
			\end{figure}
		}% Original if not 'y'
	}
	{
		\ifthenelse{\equal{#2}{w}}
		{
			\PackageWarning{bmstu}{The 'wrapfigure' environment and rotation ('y') are not simultaneously implemented in this example. Using unrotated version.}
			\begin{wrapfigure}{#3}{#4}
				\centering
				\includegraphics[width=#4]{inc/img/#1}
				\caption{#5}
				\label{img:#1}
			\end{wrapfigure}
		}
		{
			\PackageError{bmstu}{unknown image type}{}    
		}
	}
}


% Переопределение списков
\setlist[itemize]{label=--}